\subsection{Cross-cohort evaluation}
\label{sec:evaluation_design}
Each LOCO round held out one analytical cohort and used the other three for model development (Table~\ref{tab:splits}). Saved sample assignments confirmed no train/test overlap within a round. Some cohorts comprise multiple BioProjects, so this is cohort-level rather than six-study-level validation.

\begin{table}[!t]
\caption{Outer LOCO assignments. Each held-out cohort is never used for its own gene selection or model training.}
\label{tab:splits}
\centering\footnotesize
\begin{tabular}{@{}lrrl@{}}
\toprule
Test cohort & Train $n$ & Test $n$ & Development cohorts\\
\midrule
D1 & 429 & 260 & D2, D3, D4\\
D2 & 510 & 179 & D1, D3, D4\\
D3 & 613 & 76 & D1, D2, D4\\
D4 & 515 & 174 & D1, D2, D3\\
\bottomrule
\end{tabular}
\end{table}

LASSO and Elastic Net tuned regularisation using three inner cohort-held-out validation splits within each development set, selecting settings by mean validation AUROC. Other methods retained their method-specific selection and tuning procedures. In the sparse logistic pipelines, expression was transformed to $\log_2(1+\mathrm{CPM})$, where CPM used counts summed over the shared genes in each sample. Genes with CPM $>1$ in at least 10\% of the fitting samples and nonzero variance were retained. Gene filtering and standardisation were fitted only on each relevant training partition and then applied to validation or test data. Native methods used their own documented preprocessing. Method-specific prediction outputs were retained separately and are not treated as matched classifier comparisons.

\subsection{Controlled downstream classifier}
\label{sec:common_classifier}
To compare the predictive usefulness of selected genes under the same downstream model, we used each method's frozen, outer-training-selected panel in a separate L2-regularised logistic regression. Panels comprised the nonzero sparse-model coefficients, Stability Selection genes at a 0.60 frequency cutoff, native StableMate stable-significant genes, and the TENTACLES consensus. All panels were evaluated on the same four LOCO splits, with $\log_2(1+\mathrm{CPM})$ expression, training-only standardisation and class-balanced fitting (scikit-learn \texttt{liblinear}, $C=1$, 5,000 maximum iterations, seed 123). CD was positive; probabilities of at least 0.5 were classified as CD. The learner and evaluation settings were fixed, but the upstream feature selection, tuning and gene counts differed.

\subsection{Prediction measures}
\label{sec:predictive_metrics}
CD was the positive class for the verified Python prediction pipelines. We report AUROC and balanced accuracy (mean sensitivity and specificity), which helps describe performance despite differing CD prevalence. Threshold-based measures used a CD probability threshold of 0.5; F1-score and average precision are reported where available. Each cohort's result is shown separately, with an unweighted four-cohort mean and, where available, the standard deviation across cohorts.

Within-cohort stratified bootstrap intervals for AUROC used 1,000 resamples for LASSO/Elastic Net and 2,000 for StableMate. These intervals reflect test-sample variation conditional on the fitted models; they do not estimate uncertainty over future independent cohorts.

\subsection{Gene and pathway agreement}
\label{sec:reproducibility_metrics}
Let $S_{m,f}$ denote the genes selected by method $m$ when cohort $f$ is held out. For each method, we calculated the Jaccard similarity between all six unordered pairs of its four LOCO signatures and averaged them:
\begin{equation}
\bar{J}_m=\frac{1}{6}\sum_{a<b}
\frac{|S_{m,a}\cap S_{m,b}|}{|S_{m,a}\cup S_{m,b}|}.
\label{eq:jaccard}
\end{equation}
We also report panel sizes and the number of genes appearing in all four panels. LOCO training sets overlap in two cohorts, so the six pairwise scores are dependent measures of \emph{selection consistency}, not six independent replications. Signatures selected separately within single cohorts provide a distinct comparison where available.

Because Jaccard similarity depends partly on gene-set size, we additionally compared the four highest-ranked genes within each LOCO panel for LASSO, Elastic Net, Stability Selection and StableMate. Rankings used absolute standardised coefficients, subsample selection frequency, and native predictive importance, respectively; ties were settled by gene name. Four genes were the largest common panel size available across these four methods and folds. An exploratory top-ten comparison covered only LASSO, Elastic Net and StableMate. TENTACLES V2 was excluded from the fixed top-four analysis: its genuine vote counts provide coarse evidence of importance but leave many ties among six- to nine-vote genes, so a unique meaningful within-panel top-four ordering is not available. These truncated-panel comparisons assess selection overlap, not newly fitted classifiers.

For the five methods' LOCO-selected panels, Reactome v98 pathways \cite{ragueneau2026reactome} were tested for over-representation against the eligible training-gene background. We used one-sided hypergeometric tests for pathways containing 10--500 eligible genes, followed by Benjamini--Hochberg correction ($q<0.05$) \cite{benjamini1995fdr}. Pathway Jaccard similarity compared sets of significant pathways; it was undefined when both sets were empty. Pathway overlap remains descriptive because the discovery sets overlap, and enrichment power depends on panel size and gene annotation coverage.
